% Frontera de truncamiento de la vía dodecafásica completa y delimitación
% exacta del resolvente analítico finito.
\paragraph{Dépendance exacte à l'égard de la frontière de troncature.}
\label{ampl-alfa-frontera-exacta}
Les coordonnées régionales et le registre dodécaphasique procèdent de la
génération commune APP--TRIT--TPK. Une fois produites, leur composition
positionnelle possède une dépendance explicite à l'égard de la profondeur
de lecture. Soient $B=1000$,
\[
 Z_j=P_j+E_j-\Phi_j-K_j^{\mathrm{per}},\qquad
 y_N=\sum_{j=1}^{N}Z_jB^{-j},\qquad
 R_{N+1}(Z)=\sum_{r\geq0}Z_{N+1+r}B^{-r-1}.
\]
La série complète de \eqref{eq:alpha-serie} satisfait exactement
\[
 \alpha_A=y_N+B^{-N}R_{N+1}(Z),\qquad
 -2<R_{N+1}(Z)<2.
\]
Cette identité réunit la contribution de la fenêtre et celle de son
prolongement, toutes deux exprimées dans les mêmes coordonnées.

\begin{proposition}[Transport de la frontière et cylindre canonique]
\label{ampl-alfa-prop-frontera}
Soit $A_j^{(N,t)}$ la sortie de la division positionnelle finie avec
$c_{N+1}=t\in\mathcal C=\{-2,-1,0,1,2\}$, et soit
$a_N^{(t)}=\sum_{j=1}^{N}A_j^{(N,t)}B^{-j}$. Pour les registres du
domaine considéré, on a
\begin{equation}
 \alpha_A-a_N^{(t)}
 =B^{-N}\bigl(R_{N+1}(Z)-t\bigr),
 \qquad
 |\alpha_A-a_N^{(t)}|<4B^{-N}.
 \label{ampl-alfa-eq-frontera}
\end{equation}
Le quotient du prolongement complet est
$c_{N+1}^{\infty}=\lfloor R_{N+1}(Z)\rfloor$. Avec cette frontière,
\[
 0\leq\alpha_A-a_N^{(c_{N+1}^{\infty})}<B^{-N},
\]
de sorte que la fenêtre obtenue est le préfixe canonique d'ordre $N$.
\end{proposition}
\begin{proof}
Le télescopage de \eqref{eq:alpha-identidad-local} donne
$a_N^{(t)}=y_N+tB^{-N}-c_1$. Comme $Z_1=7$, chaque état entrant
appartient à $\mathcal C$ et le premier dividende est compris entre $5$ et $9$,
on a $c_1=0$. La substitution de cette égalité dans la décomposition de la
série prouve \eqref{ampl-alfa-eq-frontera}. La majoration résulte de
$R_{N+1}(Z)\in(-2,2)$ et de $t\in\mathcal C$.

Pour la normalisation canonique complète,
$c_{N+1}^{\infty}=R_{N+1}(Z)-R_{N+1}(A)$ et
$R_{N+1}(A)\in[0,1)$. Le caractère entier du quotient implique
$c_{N+1}^{\infty}=\lfloor R_{N+1}(Z)\rfloor$. La formule d'erreur
devient alors
$\alpha_A-a_N^{(c_{N+1}^{\infty})}=B^{-N}R_{N+1}(A)$, ce qui prouve
l'appartenance au cylindre canonique.
\end{proof}

La proposition identifie la signification de la frontière droite : c'est la
partie entière de la queue signée exprimée à l'échelle de la position
suivante. Le calcul sur les cinq frontières conserve cette information
comme une collection finie de possibilités jusqu'à ce que les blocs de
gauche se stabilisent. La contraction quantitative provient de la
profondeur et du transport de la queue ; son objet est la normalisation des
coordonnées déjà engendrées.

\paragraph{Domaine gradué et statut de la représentation analytique.}
\label{ampl-alfa-dominio-resolvente}
La composition analytique présentée dans \S\ref{subsec:alpha-resolvente} agit sur $\operatorname{span}\{e_7,e_8,e_9\}$ et satisfait $Ne_9=0$. Son évaluation produit exactement $T_{7:9}$ ; le degré $9$ est la dernière composante de ce domaine. L'identité $N^3=0$ explique pourquoi la résolvante s'exprime au moyen de trois termes. La résolvante locale ne détermine pas à elle seule les coefficients de degré supérieur. La section~\ref{subsec:alpha-limite-dos-vias} démontre cette non-unicité du jet nu et distingue deux opérations : la représentation dodécaphasique produit la coordonnée \(\alpha_{\mathrm{HMT}}\) à partir des registres de l'état ; ensuite, la représentation analytique utilise la même précoordonnée et engendre, par une récurrence explicite, tous les coefficients de son image analytique. Aucune racine conventionnelle n'est consultée et la représentation analytique n'est pas promue au rang de générateur indépendant. Les carrés de compatibilité et les cylindres emboîtés prouvent l'égalité des deux représentations à chaque profondeur. La nilpotence locale, la représentation dodécaphasique et la représentation analytique remplissent ainsi des fonctions distinctes : la première certifie le résultat fini d'ordre neuf ; la deuxième détermine \(\alpha_{\mathrm{HMT}}\) ; la troisième représente cette même coordonnée à toute profondeur et en vérifie la compatibilité.
